% ============================================================================
% Bab 13: Optimisasi Multiobjektif
% Sumber buku: part1-linear-programming/ch09-multi-objective/
%              multi-objective-optimization_updated.tex
% Penomoran latihan diselaraskan dengan urutan \begin{ex} per 2026-07-13.
% Semua perhitungan diverifikasi secara numerik (PuLP 3.3.2/CBC untuk LP).
% ============================================================================

\section{Bab 13: Optimisasi Multiobjektif}

Bab ini memiliki sepuluh latihan. Bagian pemanasan (13.1--13.3) mencakup
definisi optimalitas Pareto, pemeriksaan dominasi dalam sebuah tabel kecil, dan
penilaian jumlah berbobot. Soal inti (13.4--13.8) menerapkan metode-metode utama
bab ini pada data konkret: identifikasi batas, jumlah berbobot, metode
kendala-$\varepsilon$ (pada sebuah tabel, pada dua LP kecil, dan secara manual),
serta analisis kompromi. Soal konsep (13.9--13.10, dengan 13.10 sebagai soal
tantangan) menyelidiki titik buta jumlah berbobot yang telah dikenal pada batas
tak konveks dan menghitung interval bobot tepat yang memilih setiap titik sudut
dalam contoh furnitur bab ini. Setiap pernyataan numerik di bawah diverifikasi
dengan perhitungan langsung; latihan berbasis LP (13.6 dan penyapuan furnitur)
juga diselesaikan ulang dengan PuLP/CBC, dan keluaran yang ditampilkan adalah
keluaran aktual. Solusi Terpilih buku untuk 13.1, 13.4, 13.7, dan 13.10 telah
diperiksa dan terbukti benar; solusi-solusi itu digunakan kembali dan diperluas
di sini.

% ----------------------------------------------------------------------------
\exsol{13.1}{Mendefinisikan Optimalitas Pareto}{1}

Untuk masalah yang meminimumkan $f_1$ dan $f_2$ pada himpunan layak $P$, sebuah
titik $\x \in P$ bersifat \emph{optimal Pareto} jika tidak ada
$\overline{\x} \in P$ dengan $f_i(\overline{\x}) \leq f_i(\x)$ untuk kedua
$i = 1, 2$ dan $f_i(\overline{\x}) < f_i(\x)$ untuk sedikitnya satu $i$.
Dengan kata lain: tidak ada titik layak lain yang sedikitnya sama baik dalam
kedua tujuan dan secara tegas lebih baik dalam salah satunya. (Secara ekuivalen,
sesuaikan definisi maksimisasi buku dengan membalik pertidaksamaan.) Batas
Pareto adalah himpunan semua titik tersebut.

% ----------------------------------------------------------------------------
\exsol{13.2}{Menemukan Desain Terdominasi}{1}

Datanya (kedua tujuan diminimumkan):
$G = (2,9)$, $H = (4,7)$, $I = (5,8)$, $J = (7,3)$, $K = (6,5)$.

\begin{enumerate}
\item Hanya kontrak $I$ yang terdominasi: $H = (4,7)$ mengalahkan
$I = (5,8)$ dalam biaya maupun waktu pengiriman. Dengan memeriksa setiap
pasangan lain, satu kontrak selalu lebih murah sedangkan yang lain lebih cepat,
sehingga tidak ada dominasi lain (misalnya, $K = (6,5)$ versus $J = (7,3)$:
$K$ lebih murah, $J$ lebih cepat).
\item Menurut urutan $f_1$ yang meningkat, batas Paretonya adalah
\[
G = (2,9), \quad H = (4,7), \quad K = (6,5), \quad J = (7,3).
\]
Perhatikan bahwa nilai $f_2$ turun secara tegas di sepanjang daftar ini,
sebagaimana mestinya pada batas masalah minimisasi dua tujuan.
\item Benar. Jika $G$ adalah peminimum \emph{unik} $f_1$, setiap kontrak lain
memiliki $f_1$ yang secara tegas lebih besar, sehingga tidak ada kontrak yang
dapat sedikitnya sama baik dengan $G$ dalam kedua tujuan; oleh karena itu,
tidak ada yang mendominasi $G$. (Keunikan penting: jika dua kontrak memiliki
nilai minimum $f_1$ yang sama, kontrak dengan $f_2$ lebih besar terdominasi
oleh kontrak lainnya.)
\end{enumerate}

% ----------------------------------------------------------------------------
\exsol{13.3}{Menghitung Skor Jumlah Berbobot}{1}

\begin{enumerate}
\item Dengan $(w_1, w_2) = (0.6, 0.4)$, $S = 0.6 f_1 + 0.4 f_2$:
\begin{center}
\begin{tabular}{c|ccccc}
Kontrak & G & H & I & J & K \\\hline
$S$ & $\mathbf{4.8}$ & $5.2$ & $6.2$ & $5.4$ & $5.6$ \\
\end{tabular}
\end{center}
Sebagai contoh, $S_G = 0.6(2) + 0.4(9) = 1.2 + 3.6 = 4.8$. Pemenangnya adalah
kontrak $G$.
\item Dengan $(w_1, w_2) = (0.2, 0.8)$:
\begin{center}
\begin{tabular}{c|ccccc}
Kontrak & G & H & I & J & K \\\hline
$S$ & $7.6$ & $6.4$ & $7.4$ & $\mathbf{3.8}$ & $5.2$ \\
\end{tabular}
\end{center}
Pemenangnya kini kontrak $J$. Pemenang berubah karena bobot menyandikan nilai
tukar pengambil keputusan di antara tujuan-tujuan: dengan sebagian besar bobot
pada waktu pengiriman ($w_2 = 0.8$), kontrak dengan waktu pengiriman yang jauh
terbaik ($J$, yaitu 3 hari) mengatasi biayanya yang tinggi, sedangkan dengan
bobot yang sebagian besar pada biaya, kontrak termurah ($G$) menang. Secara
geometris, kedua vektor bobot bersesuaian dengan garis pendukung yang memiliki
kemiringan berbeda dan menyentuh batas di titik-titik yang berbeda.
\item $S_I - S_H = w_1(5 - 4) + w_2(8 - 7) = w_1 + w_2 > 0$ setiap kali
$w_1, w_2 > 0$. Jadi, skor $I$ secara tegas lebih buruk daripada $H$ untuk
setiap pasangan bobot positif, dan secara umum sebuah titik terdominasi kalah
dari setiap titik yang mendominasinya: selisih skor merupakan kombinasi tak
negatif dari selisih per koordinat, dengan sedikitnya satu selisih yang secara
tegas positif.
\end{enumerate}

% ----------------------------------------------------------------------------
\exsol{13.4}{Batas Pareto, Jumlah Berbobot, dan Kendala-$\varepsilon$}{2}

Kandidatnya (kedua tujuan diminimumkan):
$A = (6,5)$, $B = (4,8)$, $C = (7,4)$, $D = (3,6)$, $E = (5,6)$,
$F = (10,3)$.

\begin{enumerate}
\item Bandingkan pasangan untuk menentukan dominasi. $D = (3,6)$ mendominasi
$B = (4,8)$, dan $E = (5,6)$ terdominasi oleh $D$ ($3 \leq 5$, $6 \leq 6$,
dengan pertidaksamaan pertama tegas). $A$, $C$, dan $F$ tidak terdominasi oleh
titik mana pun. Batas Paretonya adalah $\{D, A, C, F\}$ dengan vektor tujuan
$(3,6)$, $(6,5)$, $(7,4)$, $(10,3)$.
\item Sebuah desain menjadi optimum jumlah berbobot untuk suatu $\alpha$ tepat
ketika desain tersebut berada pada selubung konveks kiri bawah dari titik-titik
batas. Dengan menggambarkan $(3,6)$, $(6,5)$, $(7,4)$, $(10,3)$: selubung dibentuk
oleh $D$, $C$, dan $F$, sedangkan $A$ berada di atasnya. Memang, $D$ ($f_1$
terbaik) dan $F$ ($f_2$ terbaik) selalu dapat diperoleh pada bobot ekstrem;
$C$ berada pada selubung di antara keduanya (ruas $D \to F$ memiliki kemiringan
$-3/7$, dan pada $f_1 = 7$ ruas itu berada pada
$6 - 4 \cdot \tfrac{3}{7} \approx 4.29 > 4$, sehingga $C$ berada di bawah ruas
itu dan karenanya pada selubung); tetapi $A = (6,5)$ berada di atas ruas dari
$D$ ke $C$ (ruas tersebut pada $f_1 = 6$ bernilai
$6 - \tfrac{2}{4}(6-3) = 4.5 < 5$). Jadi, $D$, $C$, dan $F$ dapat diperoleh
dengan jumlah berbobot, sedangkan $A$ tidak.
\item Pilih $\varepsilon = 5$: meminimumkan $f_1$ dengan syarat
$f_2 \leq 5$ menyisakan kandidat $A = (6,5)$, $C = (7,4)$, dan $F = (10,3)$,
serta memilih $A = (6,5)$, tepat titik batas yang tidak dapat dicapai dengan
jumlah berbobot mana pun. (Setiap $\varepsilon \in [5, 6)$ berhasil.)
\item Deskripsi sketsa: gambarkan keenam titik pada bidang $(f_1, f_2)$. Tandai $B$
dan $E$ sebagai terdominasi ($B$ berada di atas dan di kanan $D$; $E$ berada
tepat di kanan $D$). Batasnya adalah ``tangga'' titik-titik $D, A, C, F$ yang
membentang dari kiri atas ke kanan bawah; anak panah perbaikan mengarah ke kiri
(biaya lebih rendah) dan ke bawah (emisi lebih rendah), menuju wilayah kosong
di kiri bawah. Sketsa yang benar menunjukkan bahwa tidak ada titik batas yang
memiliki titik lain di kuadran kiri bawahnya.
\end{enumerate}

% ----------------------------------------------------------------------------
\exsol{13.5}{Metode Jumlah Berbobot}{2}

Dalam bentuk minimisasi: $P1 = (f_1, f_2) = (2, -7)$, $P2 = (5, -9)$,
$P3 = (3, -6)$.

\begin{enumerate}
\item $P1$ dan $P2$ optimal Pareto; $P3$ tidak. $P1$ mendominasi $P3$: biayanya
lebih rendah ($2 < 3$) \emph{dan} skor aksesibilitasnya lebih tinggi ($7 > 6$).
Di antara $P1$ dan $P2$ terdapat kompromi yang nyata ($P1$ lebih murah,
$P2$ lebih mudah diakses), sehingga keduanya tidak saling mendominasi.
\item Skor $S = 0.5 f_1 + 0.5 f_2$:
\begin{align*}
S_{P1} &= 0.5(2) + 0.5(-7) = -2.5, \\
S_{P2} &= 0.5(5) + 0.5(-9) = -2.0, \\
S_{P3} &= 0.5(3) + 0.5(-6) = -1.5.
\end{align*}
Skor terendah menang: proyek $P1$ dipilih.
\item Tidak. Karena $P1$ mendominasi $P3$,
$S_{P3} - S_{P1} = w_1(3-2) + w_2(-6+7) = w_1 + w_2 > 0$ untuk semua
$w_1, w_2 > 0$, sehingga skor $P3$ selalu secara tegas lebih buruk daripada
$P1$. Bahkan bobot batas degenerat gagal memilih $P3$: dengan $w = (1,0)$
(hanya biaya), $P1$ menang, dan dengan $w = (0,1)$ (hanya aksesibilitas), $P2$
menang. Alternatif terdominasi tidak pernah dapat dipilih dengan metode jumlah
berbobot yang menggunakan bobot positif.
\end{enumerate}

% ----------------------------------------------------------------------------
\exsol{13.6}{Metode Kendala-Epsilon}{2}

\begin{enumerate}
\item Dengan mempertahankan $f_1$ sebagai tujuan dan membatasi $f_2$:
\begin{align*}
\min \quad & x_1 + 2x_2 \\
\st & 3x_1 + x_2 \leq \varepsilon, \\
& x_1 + x_2 \geq 4, \\
& 0 \leq x_1 \leq 5, \quad 0 \leq x_2 \leq 5.
\end{align*}
\item Untuk $\varepsilon = 10$: optimumnya adalah $(x_1, x_2) = (3, 1)$ dengan
$f_1 = 5$ dan $f_2 = 10$. Untuk $\varepsilon = 6$: optimumnya adalah $(1, 3)$
dengan $f_1 = 7$ dan $f_2 = 6$. Dalam kedua kasus, optimum merupakan
perpotongan $3x_1 + x_2 = \varepsilon$ dengan $x_1 + x_2 = 4$ (di sepanjang
$x_1 + x_2 = 4$, kita memiliki $f_1 = 8 - x_1$, sehingga LP mendorong $x_1$
setinggi yang diizinkan batas $\varepsilon$: $2x_1 = \varepsilon - 4$). Kedua
solusi optimal Pareto: masing-masing adalah peminimum unik $f_1$ dengan syarat
$f_2 \leq \varepsilon$, dan batas $\varepsilon$ aktif, sehingga tidak ada titik
layak yang memperbaiki salah satu tujuan tanpa memperburuk yang lain.
Telah diverifikasi dengan PuLP:

\begin{verbatim}
import pulp

for eps in [10, 6]:
    prob = pulp.LpProblem("biobj", pulp.LpMinimize)
    x1 = pulp.LpVariable("x1", lowBound=0, upBound=5)
    x2 = pulp.LpVariable("x2", lowBound=0, upBound=5)
    prob += x1 + 2*x2                 # f1 (tujuan utama)
    prob += x1 + x2 >= 4
    prob += 3*x1 + x2 <= eps          # f2 <= epsilon
    prob.solve(pulp.PULP_CBC_CMD(msg=0))
    print(f"eps={eps}: x=({x1.value()}, {x2.value()}), "
          f"f1={x1.value() + 2*x2.value()}, "
          f"f2={3*x1.value() + x2.value()}")
\end{verbatim}

Keluaran aktual:
\begin{verbatim}
eps=10: x=(3.0, 1.0), f1=5.0, f2=10.0
eps=6: x=(1.0, 3.0), f1=7.0, f2=6.0
\end{verbatim}

\item Meminimumkan $f_2$ saja menghasilkan $f_2^* = 4$ pada $(0, 4)$;
meminimumkan $f_1$ saja menghasilkan $f_1 = 4$ pada $(4, 0)$, tempat
$f_2 = 12$. Jadi: untuk $\varepsilon \geq 12$, batas tidak aktif dan solusi
tetap pada $(4,0)$ dengan $(f_1, f_2) = (4, 12)$. Ketika $\varepsilon$ menurun
dari $12$ menuju $f_2^* = 4$, optimum bergeser di sepanjang sisi
$x_1 + x_2 = 4$ dari $(4,0)$ ke $(0,4)$, dengan
$x_1 = (\varepsilon - 4)/2$ dan
$f_1 = 8 - (\varepsilon - 4)/2$ yang meningkat secara linear dari $4$ ke $8$:
setiap unit pengetatan $f_2$ mengorbankan tepat $\tfrac12$ unit $f_1$.
Penyapuan menelusuri seluruh batas Pareto, yaitu ruas dari $(4,12)$ ke $(8,4)$
di ruang tujuan. Untuk $\varepsilon < 4$, LP tidak layak, sehingga penyapuan
berhenti pada $f_2^*$.
\end{enumerate}

% ----------------------------------------------------------------------------
\exsol{13.7}{Kendala-Epsilon secara Manual}{2}

\begin{enumerate}
\item Pada $(0,0)$: $(f_1, f_2) = (0,0)$. Pada $(6,0)$: $(18, 6)$. Pada
$(0,6)$: $(6, 18)$. Jadi, $(6,0)$ memaksimumkan $f_1$ saja dan $(0,6)$
memaksimumkan $f_2$ saja.
\item Untuk $\varepsilon > 6$, kendala $x_1 + 3x_2 \geq \varepsilon$ aktif
pada optimum, dan solusinya berada di perpotongan kendala itu dengan sisi
$x_1 + x_2 = 6$: penyelesaian kedua persamaan menghasilkan
$x_2 = (\varepsilon - 6)/2$ dan $x_1 = (18 - \varepsilon)/2$. (Untuk
$\varepsilon = 6$, titik sudut $(6,0)$ sendiri optimal.) Dalam bentuk tabel:
\begin{center}
\begin{tabular}{c|cccc}
$\varepsilon$ & $x_1$ & $x_2$ & $f_1$ & $f_2$ \\\hline
6  & 6 & 0 & 18 & 6 \\
10 & 4 & 2 & 14 & 10 \\
14 & 2 & 4 & 10 & 14 \\
18 & 0 & 6 & 6  & 18 \\
\end{tabular}
\end{center}
Setiap baris mudah dikonfirmasi secara grafis: wilayah yang diiris adalah
segitiga yang dipotong oleh garis $x_1 + 3x_2 = \varepsilon$, dan tujuan
$3x_1 + x_2$ meningkat ke arah kanan bawah, sehingga optimum berada di sudut
irisan pada hipotenusa.
\item Setiap solusi yang ditabulasikan memiliki tepat $f_2 = \varepsilon$ dan
$f_1$ terbesar yang dimungkinkan dengan batas tersebut, sehingga setiap solusi
optimal Pareto. Batas Pareto di ruang keputusan adalah sisi $x_1 + x_2 = 6$
dari $(6,0)$ ke $(0,6)$: parameterisasi dengan $x_2 = t$ menghasilkan
$f_1 = 18 - 2t$ dan $f_2 = 6 + 2t$, sehingga laju komprominya satu banding satu
($\Delta f_1 / \Delta f_2 = -1$): setiap unit yang diperoleh dalam $f_2$
mengorbankan tepat satu unit $f_1$.
\item Nilai maksimum $f_2$ di atas $X$ adalah $18$ (pada $(0,6)$), sehingga
untuk $\varepsilon > 18$, LP berkendala tersebut \emph{tidak layak}; penyapuan
harus berhenti pada $\varepsilon = 18$. (Pelajaran praktis bagi implementasi:
pertama-tama hitung optimum individual setiap tujuan untuk menetapkan rentang
penyapuan.)
\end{enumerate}

% ----------------------------------------------------------------------------
\exsol{13.8}{Analisis Kompromi}{2}

\begin{enumerate}
\item Deskripsi grafik: keempat titik $(40,20)$, $(45,14)$, $(55,10)$, $(70,8)$
pada bidang (biaya, emisi) membentuk rantai menurun; menghubungkannya dari kiri
atas ke kanan bawah menghasilkan batas Pareto (diskret dan konveks). Tidak ada
titik yang berada di kiri bawah titik lain, sehingga keempatnya berada pada
batas; arah perbaikan adalah menuju kiri bawah.
\item Tingkat substitusi marginal (emisi yang dikurangi per \$1000 biaya
tambahan):
\begin{center}
\begin{tabular}{l|ccc}
Langkah & $\Delta f_1$ & $\Delta f_2$ & ton per \$1000 \\\hline
S1 $\to$ S2 & $+5$  & $-6$ & $6/5 = 1.20$ \\
S2 $\to$ S3 & $+10$ & $-4$ & $4/10 = 0.40$ \\
S3 $\to$ S4 & $+15$ & $-2$ & $2/15 \approx 0.13$ \\
\end{tabular}
\end{center}
Langkah S1 $\to$ S2 memberikan pengurangan emisi per dolar yang jauh terbesar
(1.2 ton per \$1000). Rasio yang terus memburuk ($1.20 \to 0.40 \to 0.13$)
berarti S2 adalah siku (titik lutut) alami batas ini.
\item Dengan anggaran paling banyak $\$50{,}000$, solusi Pareto yang layak adalah
S1 ($\$40$ ribu) dan S2 ($\$45$ ribu); perusahaan sebaiknya memilih \textbf{S2}:
solusi itu sesuai anggaran dan memiliki emisi yang lebih rendah (14 versus
20 ton) daripada S1, sedangkan S3 ($\$55$ ribu) tidak terjangkau. (Selaras dengan
bagian (2), S2 juga kebetulan menjadi langkah dengan nilai terbaik pada batas.)
\end{enumerate}

% ----------------------------------------------------------------------------
\exsol{13.9}{Hal yang Tidak Dapat Dilihat oleh Jumlah Berbobot}{2}

Titik-titiknya, dengan kedua tujuan diminimumkan: $P = (0,4)$, $Q = (3,3)$,
$R = (4,0)$.

\begin{enumerate}
\item Secara berpasangan: $P$ versus $Q$: $P$ lebih baik dalam $f_1$ ($0 < 3$),
$Q$ lebih baik dalam $f_2$ ($3 < 4$). $Q$ versus $R$: $Q$ lebih baik dalam
$f_1$, $R$ lebih baik dalam $f_2$. $P$ versus $R$: $P$ lebih baik dalam $f_1$,
$R$ lebih baik dalam $f_2$. Tidak ada dominasi pada pasangan mana pun, sehingga
ketiga titik optimal Pareto.
\item Skor berbobot dengan bobot $(\alpha, 1-\alpha)$:
\[
S_P(\alpha) = 4(1-\alpha) = 4 - 4\alpha, \qquad
S_Q(\alpha) = 3\alpha + 3(1-\alpha) = 3, \qquad
S_R(\alpha) = 4\alpha.
\]
Untuk setiap $\alpha \in [0,1]$,
$\min\{S_P, S_R\} = \min\{4 - 4\alpha, \, 4\alpha\} \leq 2$, karena kedua
nilai berjumlah $4$ (nilai minimum keduanya dimaksimumkan pada
$\alpha = \tfrac12$, saat keduanya sama dengan $2$). Karena itu,
$\min\{S_P, S_R\} \leq 2 < 3 = S_Q$ untuk setiap $\alpha$: $P$ atau $R$ selalu
memiliki skor yang secara tegas lebih baik daripada $Q$, sehingga $Q$ tidak
pernah menjadi optimum jumlah berbobot, bahkan sebagai nilai seri.
\item Secara geometris, metode jumlah berbobot meminimumkan fungsi linear di
atas titik-titik di ruang tujuan, sehingga optimumnya tepat merupakan titik
pada selubung konveks kiri bawah, tempat suatu garis pendukung dengan normal
$(\alpha, 1-\alpha) \geq 0$ menyentuh. Di sini batas bawah selubung adalah ruas
dari $P = (0,4)$ ke $R = (4,0)$; titik tengahnya adalah $(2,2)$, dan
$Q = (3,3)$ berada secara tegas di atas ruas tersebut (garis yang melalui $P$
dan $R$ adalah $f_1 + f_2 = 4$, sedangkan $Q$ memiliki $f_1 + f_2 = 6$).
Setiap garis pendukung hanya menyentuh $P$, $R$, atau seluruh ruas, tidak pernah
$Q$. Inilah tepatnya entri ``Kekurangan'' untuk metode jumlah berbobot pada
Tabel~13.1: metode ini dapat melewatkan solusi optimal Pareto di wilayah tak
konveks pada batas, dan batas tiga titik ini tak konveks pada $Q$.
\item Metode kendala-$\varepsilon$ memperoleh kembali $Q$: minimumkan $f_1$
dengan syarat $f_2 \leq \varepsilon$ menggunakan $\varepsilon = 3$. Titik-titik
yang layak kemudian adalah $Q$ (dengan $f_2 = 3$) dan $R$ (dengan $f_2 = 0$),
dan $Q$ menang karena $f_1 = 3 < 4$. (Pemrograman sasaran dengan sasaran $(3,3)$
dan sembarang bobot positif juga memilih $Q$, yang memiliki simpangan total $0$.)
\end{enumerate}

% ----------------------------------------------------------------------------
\exsol{13.10}{Bobot yang Memilih Setiap Titik Sudut}{3}

Solusi Terpilih buku benar sebagaimana tercetak dan diperluas di sini;
semua titik pisah di bawah diverifikasi ulang baik secara aljabar maupun dengan
penyapuan PuLP terhadap $\alpha$.

\begin{enumerate}
\item Tujuan berbobot
$\alpha(8x + 2y) - (1-\alpha)(10x + 2y)$ pada setiap titik sudut $P$:
\[
(0,0): \; 0, \qquad
(0,20): \; 80\alpha - 40, \qquad
(10,8): \; 212\alpha - 116, \qquad
(12,0): \; 216\alpha - 120.
\]
(Sebagai contoh, $(10,8)$ memiliki $f_1 = 96$ dan $f_2 = 116$, yang memberikan
$96\alpha - 116(1 - \alpha) = 212\alpha - 116$.)
\item Nilai optimal merupakan selubung atas dari keempat fungsi linear
$\alpha$ ini. Dengan membandingkan pasangan: $(0,20)$ mengalahkan $(0,0)$ jika
dan hanya jika $80\alpha - 40 \geq 0$, yaitu $\alpha \geq \tfrac12$; $(10,8)$
mengalahkan $(0,20)$ jika dan hanya jika
$212\alpha - 116 \geq 80\alpha - 40$, yaitu $132\alpha \geq 76$, yakni
$\alpha \geq \tfrac{19}{33} \approx 0.576$; dan $(10,8)$ mengalahkan $(12,0)$
jika dan hanya jika $4 \geq 4\alpha$, yang berlaku untuk semua
$\alpha \leq 1$. Jadi,
\[
(0,0) \text{ optimal untuk } \alpha \in [0, \tfrac12], \qquad
(0,20) \text{ untuk } \alpha \in [\tfrac12, \tfrac{19}{33}], \qquad
(10,8) \text{ untuk } \alpha \in [\tfrac{19}{33}, 1],
\]
dengan titik pisah $\alpha = \tfrac12$ dan $\alpha = \tfrac{19}{33}$.
Penyapuan numerik mengonfirmasi peralihan tersebut (keluaran aktual PuLP/CBC
dari LP berbobot): $\alpha = 0.4999$ menghasilkan $(0,0)$;
$\alpha = 0.5001$ dan $\alpha = 0.5757$ menghasilkan $(0,20)$;
$\alpha = 0.5758$ menghasilkan $(10,8)$.
\item Pada $\alpha = \tfrac12$, seluruh sisi $x = 0$ optimal (setiap rencana
$(0, y)$, $0 \leq y \leq 20$, berskor $0$); pada
$\alpha = \tfrac{19}{33}$, seluruh sisi sonokeling dari $(0,20)$ hingga $(10,8)$
optimal. Pada $\alpha = 1$, limbah diabaikan, sehingga semua rencana yang
optimal menurut pendapatan bersifat optimal: seluruh sisi dari $(10,8)$ hingga
$(12,0)$, \emph{termasuk titik sudut terdominasi} $(12,0)$, yang menghasilkan
$\$96{,}000$ sama seperti $(10,8)$ tetapi membuang $120 > 116$ bdft. Pelajarannya:
jumlah berbobot dengan bobot nol dapat menghasilkan titik yang tidak optimal
Pareto; sembarang bobot yang secara tegas positif pada limbah memutus seri
untuk memenangkan $(10,8)$.
\item Dengan pendapatan dalam dolar tanpa penskalaan,
$f_1 = 8000x + 2000y$, nilai-nilai titik sudut menjadi $0$,
$40040\alpha - 40$, $96116\alpha - 116$, dan $96120\alpha - 120$, dan
perbandingan yang sama menghasilkan titik pisah
\[
\alpha = \tfrac{40}{40040} = \tfrac{1}{1001} \approx 0.001
\qquad\text{dan}\qquad
\alpha = \tfrac{76}{56076} = \tfrac{19}{14019} \approx 0.00136 :
\]
hampir setiap bobot dalam $[0,1]$ kini memilih $(10,8)$, dan seluruh struktur
kompromi dipampatkan ke dalam $0.14\%$ pertama dari rentang bobot. Ketika tujuan
berada pada skala yang sangat berbeda, bobot kehilangan makna yang dimaksudkan
sebagai kepentingan relatif, sehingga tujuan harus dinormalisasi sebelum
dibobotkan (Bagian 13.2).
\end{enumerate}
